\begin{ex}
1 chục bằng bao nhiêu đơn vị?
\shortans{10}
\loigiai{
$1\text{ chục} = 10\text{ đơn vị}$.
}
\end{ex}

\begin{ex}
10 chục bằng một trăm. Một trăm viết bằng số nào?
\shortans{100}
\loigiai{
$10\text{ chục} = 100$.
}
\end{ex}

\begin{ex}
10 trăm bằng một nghìn. Một nghìn viết bằng chữ số 1 và mấy chữ số 0?
\shortans{3}
\loigiai{
Số 1000 có 3 chữ số 0.
}
\end{ex}

\begin{ex}
Số một nghìn viết bằng chữ số nào?
\shortans{1000}
\loigiai{
Một nghìn viết là 1000.
}
\end{ex}

\begin{ex}
Số 1000 là số có mấy chữ số?
\shortans{4}
\loigiai{
Số 1000 có 4 chữ số.
}
\end{ex}

\begin{ex}
Số bé nhất có bốn chữ số là số nào?
\shortans{1000}
\loigiai{
Số bé nhất có bốn chữ số là 1000.
}
\end{ex}

\begin{ex}
Số lớn nhất có ba chữ số là số nào?
\shortans{999}
\loigiai{
Số lớn nhất có ba chữ số là 999.
}
\end{ex}

\begin{ex}
Số liền sau của số 999 là số nào?
\shortans{1000}
\loigiai{
Số liền sau của 999 là 1000.
}
\end{ex}

\begin{ex}
Số liền trước của số 1000 là số nào?
\shortans{999}
\loigiai{
Số liền trước của 1000 là 999.
}
\end{ex}

\begin{ex}
Có bao nhiêu chục trong số 100?
\shortans{10}
\loigiai{
Có 10 chục.
}
\end{ex}

\begin{ex}
Có bao nhiêu trăm trong số 1000?
\shortans{10}
\loigiai{
Có 10 trăm.
}
\end{ex}

\begin{ex}
Số gồm 3 trăm và 0 chục, 0 đơn vị là số nào?
\shortans{300}
\loigiai{
Số đó là 300.
}
\end{ex}

\begin{ex}
Số gồm 7 trăm và 0 chục, 0 đơn vị là số nào?
\shortans{700}
\loigiai{
Số đó là 700.
}
\end{ex}

\begin{ex}
Đếm thêm 100: $100, 200, 300, \dots, 500$. Số còn thiếu là số nào?
\shortans{400}
\loigiai{
Số còn thiếu là 400.
}
\end{ex}

\begin{ex}
Đếm thêm 100: $600, 700, \dots, 900, 1000$. Số còn thiếu là số nào?
\shortans{800}
\loigiai{
Số còn thiếu là 800.
}
\end{ex}

\begin{ex}
Số liền sau của số 100 là số nào?
\shortans{101}
\loigiai{
Số liền sau của 100 là 101.
}
\end{ex}

\begin{ex}
Số liền trước của số 100 là số nào?
\shortans{99}
\loigiai{
Số liền trước của 100 là 99.
}
\end{ex}

\begin{ex}
Một xưởng sản xuất đóng 10 hộp bánh, mỗi hộp có 10 chiếc bánh. Xưởng đã đóng được tất cả bao nhiêu chiếc bánh?
\shortans{100}
\loigiai{
Có: $10 \times 10 = 100$ chiếc bánh.
}
\end{ex}

\begin{ex}
Một cửa hàng có 10 thùng sách, mỗi thùng có 100 cuốn sách. Cửa hàng có tất cả bao nhiêu cuốn sách?
\shortans{1000}
\loigiai{
Có: $100 \times 10 = 1000$ cuốn sách.
}
\end{ex}

\begin{ex}
Một nghìn trừ đi một trăm bằng bao nhiêu?
\shortans{900}
\loigiai{
Ta có: $1000 - 100 = 900$.
}
\end{ex}

\begin{ex}
Số 1000 gồm mấy chữ số khác nhau? (chữ số 1 và chữ số 0). Số lượng chữ số khác nhau là bao nhiêu?
\shortans{2}
\loigiai{
Có 2 chữ số khác nhau là 1 và 0.
}
\end{ex}

\begin{ex}
Tính: $300 + 200$.
\shortans{500}
\loigiai{
Ta có $300 + 200 = 500$.
}
\end{ex}

\begin{ex}
Tính: $800 - 500$.
\shortans{300}
\loigiai{
Ta có $800 - 500 = 300$.
}
\end{ex}

\begin{ex}
Số nào cộng với 600 thì được 1000?
\shortans{400}
\loigiai{
Số đó là: $1000 - 600 = 400$.
}
\end{ex}

\begin{ex}
Số nào trừ đi 300 thì được 700?
\shortans{1000}
\loigiai{
Số đó là: $700 + 300 = 1000$.
}
\end{ex}
